\section{Ensembling: when does combining guards recover transfer?}
\label{app:ensembling}

Ensembling is a classic way to make a classifier generalize, so it is natural to ask whether it repairs
the transfer loss that fine-tuning induces (\Cref{sec:actI}). We test this \emph{retrospectively} on the
adaptation panel's committed per-row margins (\Cref{sec:adaptation}): \EnsNCk{} checkpoints across \EnsNFam{}
families (excluding the degenerate Llama-Guard null cell), each scored under \{unmodified base, SFT,
KL-SFT\} $\times$ five seeds on the \emph{same} rows. Every number is equal-family macro-AP on the raw
logit margin, anchored to each checkpoint's \emph{own} unmodified base --- identical estimand to the main
text. Within a checkpoint the base and its adapters share one output head, so raw margins are
scale-comparable and can be averaged directly; across checkpoints they are not, so the cross-model
committee (below) rank-normalizes first. This appendix is retrospective on the inspected panel, not a
preregistered claim.

\paragraph{Ensembling a fine-tune with itself is only denoising.} Averaging the five SFT seeds into one
guard raises transfer by a bootstrap-significant but small \EnsSeedGainTrans{} (one-sided LCB
\EnsSeedGainTransLCB{}) over a single SFT run --- yet transfer still lands \emph{below} the base
(\EnsSeedSFTdTrans{} vs.\ base; \Cref{tab:ensembling}). Seed-ensembling KL-SFT behaves the same
(\EnsSeedKLdTrans{}). The reason is structural: every seed is trained by the same recipe on the same data,
so all members share the \emph{same} specialization bias. Averaging them cancels seed \emph{variance}, not
the shared \emph{bias} toward the represented sources --- so it polishes a specialized guard without
un-specializing it.

\paragraph{Ensembling across the specialization axis recovers transfer.} The one ensemble that lifts
transfer above the base is the base $\oplus$ adapter combination --- which is exactly the Act~II
composition (\Cref{sec:actIII}). On this panel it raises transfer by \EnsBaseSFTvsBaseTrans{} over the base
(one-sided LCB \EnsBaseSFTvsBaseTransLCB{} $>0$), recovering \EnsBaseSFTvsSFTTrans{} relative to pure SFT,
at a represented cost of \EnsBaseSFTvsSFTRep{} (95\% CI $[\EnsBaseSFTvsSFTRepLCB, \EnsBaseSFTvsSFTRepUCB]$).
It is a recovery, not a Pareto win: for the strongest base the composed guard still gives back a little
transfer (the Act~II caveat, \Cref{sec:actIII}). \Cref{fig:ensembling-plane} makes the distinction visual:
the seed-ensembling arrows stay inside the ``transfer below base'' band, while the arrow to base $\oplus$
SFT crosses into ``transfer recovered.'' Adding KL-SFT as a third member (base $\oplus$ SFT $\oplus$
KL-SFT) lands essentially on top of base $\oplus$ SFT --- no further gain. The active ingredient is
\emph{diversity across the specialization axis}: the un-tuned base makes off-distribution errors that are
decorrelated from the tuned member's, so averaging cancels them; two specialized members do not.

% AUTO-GENERATED by experiments/emit_ensembling_tex.py -- do not hand-edit.
\begin{table}[t]\centering\small
\caption{Retrospective ensembling on the adaptation panel (\EnsNCk{} checkpoints across \EnsNFam{}
families, excl.\ the degenerate Llama-Guard null cell; equal-family macro-AP on the raw margin,
anchored to each checkpoint's \emph{own} unmodified base). Within-checkpoint ensembles average raw
margins (scale-comparable, shared head). Represented $=$ id\_test, transfer $=$ transfer\_test.
Ensembling a fine-tune \emph{with itself} (seeds) only denoises; only adding the \emph{un-tuned base}
lifts transfer above base.}
\label{tab:ensembling}
\begin{tabular}{@{}l cc cc l@{}}
\toprule
& \multicolumn{2}{c}{represented} & \multicolumn{2}{c}{transfer} & \\
\cmidrule(lr){2-3}\cmidrule(lr){4-5}
Guard / ensemble & AP & $\Delta$ & AP & $\Delta$ & effect \\
\midrule
unmodified base & \EnsBaseRep{} & --- & \EnsBaseTrans{} & --- & reference \\
\midrule
\multicolumn{6}{@{}l}{\emph{Ensemble a fine-tune with itself (redundant $\Rightarrow$ correlated errors)}} \\
SFT (single run) & $0.984$ & $+0.208$ & $0.828$ & $-0.078$ & specializes \\
SFT, 5-seed ensemble & $0.989$ & $+0.213$ & $0.852$ & $-0.053$ & denoise; still $<$ base \\
KL-SFT (single run) & $0.941$ & $+0.165$ & $0.886$ & $-0.019$ & mild specialize \\
KL-SFT, 5-seed ensemble & $0.947$ & $+0.171$ & $0.897$ & $-0.008$ & $\approx$ base \\
\midrule
\multicolumn{6}{@{}l}{\emph{Ensemble across the specialization axis (diverse $\Rightarrow$ decorrelated errors)}} \\
base $\oplus$ SFT \; \footnotesize(= Act~II)\normalsize & $0.967$ & $+0.191$ & $0.920$ & $+0.014$ & \textbf{recovers transfer} \\
base $\oplus$ KL-SFT & $0.886$ & $+0.110$ & $0.915$ & $+0.009$ & recovers, lower rep \\
SFT $\oplus$ KL-SFT & $0.986$ & $+0.210$ & $0.895$ & $-0.011$ & no transfer help \\
base $\oplus$ SFT $\oplus$ KL-SFT & $0.968$ & $+0.192$ & $0.918$ & $+0.013$ & $\approx$ base$\oplus$SFT \\
\bottomrule
\end{tabular}
\end{table}


\begin{figure}[t]\centering
\includegraphics[width=0.86\linewidth]{fig_ensembling_plane.pdf}
\caption{\textbf{The ensembling plane} (equal-family mean, \EnsNCk{}-checkpoint panel). Change in
represented (x) and held-out transfer (y) macro-AP vs.\ each checkpoint's own base. Ensembling a fine-tune
with itself (orange, single~$\to$~5 seeds) stays in the ``transfer below base'' band; only crossing the
specialization axis by adding the un-tuned base (green arrow, to base~$\oplus$~SFT) lifts transfer above
the base. \mbox{base~$\oplus$~SFT~$\oplus$~KL-SFT} coincides with base~$\oplus$~SFT (adding a second
specialized member buys nothing).}
\label{fig:ensembling-plane}
\end{figure}

\paragraph{A committee of fine-tuned guards does not generalize better.} The strongest form of the
technique --- a cross-model \emph{committee} that rank-averages different guards on the same row --- makes
the point sharpest (\Cref{tab:ensembling-committee}). A committee of the SFT guards is \emph{worse} than
the single best guard on transfer (\EnsCommSFTvsBestGen{} over the four general checkpoints,
\EnsCommSFTvsBestAll{} over all ten): SFT drives every member toward the same represented sources, so their
held-out errors are correlated and averaging cannot cancel them --- it only dilutes the best member. The
committee of \emph{KL-SFT} guards merely breaks even (\EnsCommKLvsBestGen{} / \EnsCommKLvsBestAll{}),
because KL regularization leaves the members less specialized and therefore more diverse --- but that is
the KL penalty doing the work, not the ensemble.

% AUTO-GENERATED by experiments/emit_ensembling_tex.py -- do not hand-edit.
\begin{table}[t]\centering\small
\caption{Cross-model \emph{committee} ensembles (rank-percentile average of different guards on the
same row). A committee of SFT guards does not beat the single best guard on transfer --- SFT
correlates their off-distribution errors. Only the less-specialized KL-SFT committee breaks even.}
\label{tab:ensembling-committee}
\begin{tabular}{@{}l cc c r@{}}
\toprule
Committee & rep.\ AP & transfer AP & best single (tr.) & $\Delta$ vs best \\
\midrule
\multicolumn{5}{@{}l}{\emph{4 general checkpoints}} \\
committee of bases & $0.666$ & $0.921$ & $0.942$ & $-0.022$ \\
committee of SFT guards & $0.990$ & $0.854$ & $0.868$ & $-0.014$ \\
committee of KL-SFT guards & $0.975$ & $0.918$ & $0.914$ & $+0.004$ \\
\midrule
\multicolumn{5}{@{}l}{\emph{all 10 checkpoints}} \\
committee of bases & $0.833$ & $0.970$ & $0.970$ & $-0.001$ \\
committee of SFT guards & $0.992$ & $0.876$ & $0.932$ & $-0.056$ \\
committee of KL-SFT guards & $0.975$ & $0.945$ & $0.949$ & $-0.004$ \\
\bottomrule
\end{tabular}
\end{table}


\begin{takeaway}{Does ensembling help here?}
Only when it injects the diversity fine-tuning removed. Keeping the \emph{un-tuned base} in the ensemble
(Act~II composition) recovers transfer above base; ensembling fine-tunes \emph{with each other} --- across
seeds or as a committee of tuned guards --- does not, because specialization correlates their
off-distribution errors. Ensembling is a variance tool, not a bias fix: it costs $N$ inference passes,
does not dissolve the represented/transfer tradeoff, and every composed candidate must still be
recalibrated and gated on both splits like any other (\Cref{sec:practitioner}).
\end{takeaway}
